% !TEX root = ../nc-cubic.tex
\subsection{二楕円因子を持つ五様体の無標識判定}
\label{tes:sec:fivefold}
定理\ref{csc:thm:comparison}で $k=2,m=1$ とする。実比較の標的は
$F=f_1(X)+f_2(Y)+cu^3$ ($c\ne0$) で、$E_i=V(f_i)$ は滑らかな平面三次である。
移動点、因子の順序、中心座標を残余圏に標識しない。

\begin{proposition}\label{tes:prop:unordered}
この分離族の二つの全残余 dg 圏が $\C$-線型 dg Morita 同値であることは、
重複を込めた非順序対 $\{j(E_1),j(E_2)\}$ が等しいことと同値である。
この対は内在的な二次乗法から復元できる。
さらに、滑らかな可換三次五様体との実比較が既にある圏については、
この分離族の圏に同値かどうかを次の方法で判定できる。
\end{proposition}
\begin{proof}
命題\ref{nft:prop:jacobian-necessary}で $I_2$ を取り、
全偏微分が $I_2$ を張る三次式を一つ選ぶ。
実比較があるからそのような式が存在する。
命題\ref{gcd:prop:integrability-symmetrizer}により、全ての選択は
最細直和成分をそれぞれ非零定数倍するだけで、複素数体上では
各成分の座標の三乗根倍により射影同型である。
補題\ref{ej:lem:symmetrizer-algebra}の原始冪等元でその最細分解を取る。

滑らかな平面三次が分解可能なら、一変数の三乗と平方因子を持たない
二変数三次の和に書ける。射影直線上の相異なる三点は任意の別の三点へ
射影変換で移るので、この二変数部分は二つの三乗に分解する。
従って平面三次は Fermat 型であり、逆も成立する。
これは $j=0$ と同値である。実際、複素楕円曲線の同型で移した
次数三線束の差は平行移動で補正できる（$[3]$ は全射）ため、
同じ $j$ の二つの平面モデルは射影同型である。
Fermat 平面三次の $j$ は零である。

従って最細成分の変数個数の多重集合が
$(3,3,1),(3,1,1,1,1),(1,1,1,1,1,1,1)$ のいずれかであることが判定条件である。
三変数成分から非零の $j$ を読み、残りを零で埋めて二個とする。
逆にこれらの型なら三乗を必要なだけ三個ずつまとめれば指定形の滑らかな標的を得る。
これを実現する分離表示には例えば各平面三次の通常の多項式環と方程式を
定理\ref{csc:thm:comparison}に使える。非自明な平行移動を選ぶ場合にも、
各曲線の非零三捩れ平行移動が $\mathcal O(1)$ を保つので、その射影作用を
三乗一の線型作用へ持ち上げて多項式環の Zhang twist を取れる。
方程式の指標は自明である。実際、平行移動は曲線の微分を保ち、
位数三の自由な射影作用の三固有値は相異なる（重複すると固定直線が
三次と交わる）のでこの持上げの行列式は一であり、留数式から従う。
よって三次は twist 後も中心である。
非順序対が等しければ平面モデルの射影同型とブロックの再尺度化で
標的が同型となる。この同型のグラフの核を、両端の分離比較核およびその逆と
合成すると、全残余圏間の実同値を得る。
\end{proof}

\begin{proposition}\label{tes:prop:linear-coupling}
混合稿の定理\paperref{mixed}{euk:thm:universal-kernel}で二因子、
$\mathbf t=0$ とし、標的が滑らかなファイバーを取る。
比較後の方程式を $F=G+u^2L+du^3$、$G=f_1+f_2$ と書く。
$G$ の最細分解のうち $L$ の制限が非零の成分を選ぶ。
次の規則が $I_2$ から復元される最細分解を与える。
\begin{enumerate}
\item $L=0$ なら $G$ の最細分解に一変数成分を一つ加える。
\item 選ばれた成分が一変数成分一つなら、それと $u$ の二変数部分は
二つの一変数成分に分かれる。
\item その他では、選ばれた成分の全変数と $u$ が一つの分解不能成分になり、
選ばれなかった各成分はそのまま残る。
\end{enumerate}
特に両平面三次の $j$ が非零で $L$ が両ブロックに非零なら、
全残余圏の可換標的は最細成分一個を持ち、上のどの分離表示の圏にも同値でない。
\end{proposition}
\begin{proof}
この切片で $L=\mathsf A\alpha$, $d=\kappa c$ は全ての値を動く
（混合稿の命題\paperref{mixed}{euk:prop:coefficients}）。
滑らかな点では同稿の実核により $I_2$ はこの $F$ の勾配空間である。
最初の場合は直和分解から、二番目は滑らかな二変数三次の三点から従う。

残る場合の symmetrizer を直接解く。$V$ を六変数の点の空間、
$e$ を $u$ 方向の基底、$T_G,T_F$ を偏極三重形式とし、
$T_F(e,e,x)=\ell(x)=L(x)/3$ と置く。
自己準同型を $x\mapsto Bx+h(x)e$ ($x\in V$)、$e\mapsto v+se$ と書く。
symmetrizer 条件を $V^3,V^2e,Ve^2$ で取ると、それぞれ
\[
 B\in E_G,\qquad T_G(v,x,y)=h(x)\ell(y)=h(y)\ell(x),\qquad
 \ell(Bx)+d h(x)=s\ell(x)
\]
となり、これで全条件を尽くす。$h$ は $\ell$ の定数倍である。
$\ell$ が二成分以上に非零なら $\ell^2$ は混合項を持つ一方、
$T_G(v,-,-)$ には混合項がない。従って $h=0$ である。
$\ell$ が分解不能な三変数成分一つにのみ非零の場合も、
滑らかな分解不能三次には階数一の極二次式はない。
実際 $T(v,-,-)=a\ell^2$ ($a\ne0$) なら、$\ell(v)=0$ では
$[v]$ が特異点になる。$\ell(v)\ne0$ では
$\langle v\rangle\oplus\ker\ell$ に関する全混合項が零になり分解する。
他の成分の極が零ならその成分での $v$ は零である（滑らかさ）。
こうして常に $h=0$、さらに $v=0$ を得る。
最後の条件は $B$ の各選択成分でのスカラーが $s$ であることに等しい。
$E_G$ は各最細成分への独立なスカラーからなるので、
この記述から $E_F$ の原始冪等元をすべて読める。
命題\ref{tes:prop:unordered}が無標識同値への帰結を与える。
\end{proof}

\begin{corollary}[Fermat 二因子の一次元切片]\label{tes:cor:fermat-line}
両方の境界曲線を Fermat 平面三次とし、各因子の一座標を $x,y$ とする。
比較係数で $L=x+y$ を固定し $d$ のみを動かす。
滑らかな係数は $d(16+27d^2)\ne0$ であり、この切片の全残余圏が
対応する分離型の非順序対は
\[
 \left\{0,\frac{-4096}{d^2(16+27d^2)}\right\}.
\]
従って境界二因子を固定したこの切片だけで、全ての $\{0,j\}$
($j\in\C^\times$) を実現する。
\end{corollary}
\begin{proof}
残る四座標は三乗のままであり、三つを Fermat 平面三次に、一つを中心変数に取る。
三変数部分は $x^3+y^3+u^2(x+y)+du^3$ である。
$s=x+y$, $t=x-y$ とすれば、その四倍は
$s^3+3st^2+4su^2+4du^3$ となる。$d=0$ では $s$ が因子である。
$d\ne0$ では $s=1$ の座標で $X=4du$, $Y=4d\sqrt{-3}\,t$ とすると
$Y^2=X^3+4X^2+16d^2$ を得る。この置換は射影的にも線型可逆である。
Weierstrass 不変量は $c_4=256$, $\Delta=-4096d^2(16+27d^2)$ であり、
表示した $j=c_4^3/\Delta$ と滑らかさの条件が従う。
残りの四つの三乗との分離和にも同じ滑らかさ判定が適用できる。
任意の非零有限値 $j$ に対して $j d^2(16+27d^2)=-4096$ は複素根を持ち、
各根は滑らかな係数である。命題\ref{tes:prop:unordered}と上記の普遍比較核を
合成すれば対応する全残余圏の実同値を得る。
\end{proof}
